\documentclass[11pt]{article}
\usepackage[margin=1in]{geometry}
\usepackage{amsmath,amssymb}
\usepackage{booktabs}
\usepackage{enumitem}
\usepackage{xcolor}
\usepackage{listings}
\usepackage{hyperref}
\hypersetup{colorlinks=true,linkcolor=blue!50!black,urlcolor=blue!50!black}

\lstset{basicstyle=\ttfamily\small,columns=fullflexible,breaklines=true,
  frame=single,framesep=4pt,rulecolor=\color{gray!40},backgroundcolor=\color{gray!5}}

\newcommand{\Ton}{T_{\mathrm{on}}}
\newcommand{\Toff}{T_{\mathrm{off}}}
\newcommand{\Vin}{V_{\mathrm{in}}}
\newcommand{\Vref}{V_{\mathrm{ref}}}
\newcommand{\Vout}{V_{\mathrm{out}}}
\newcommand{\Ipk}{I_{\mathrm{pk,max}}}
\newcommand{\Tsw}{T_{\mathrm{sw}}}

\title{\vspace{-1.5em}Time--Optimal $\Ton/\Toff$ Computation\\[0.3em]
\large STM32G474 Voltage--Mode Synchronous Buck --- Large--Signal Load--Step Response}
\author{G.\ Qian}
\date{\today}

\begin{document}
\maketitle
\vspace{-1em}

\begin{abstract}
\noindent This note documents exactly how the firmware computes the single
ON--OFF switching times $\Ton,\Toff$ that are applied for one stretched HRTIM
period on a detected load step--up. It covers (i) the second--order model and the
two terminal equations, (ii) their reduction to a single quadratic (one square
root) as implemented in \texttt{buck\_optimal\_onoff()}, (iii) root selection and
clamps, (iv) the ISR closed form that replaced it and its pipeline (peak--current
limit, time clamps, tick quantization), and (v) the Jacobian timing adaptation
and the post--pulse reseed. Symbols map one--to--one to the C source.\end{abstract}

\section{Plant model and shifted state}
The synchronous buck averaged over a switching action is a second--order LC plant
with a lumped series resistance $r$ (inductor DCR $+$ FET/path):
\begin{equation}
L\,\frac{di_L}{dt}=v_{\mathrm{sw}}-\Vout-r\,i_L,\qquad
C\,\frac{d\Vout}{dt}=i_L-i_{\mathrm{load}},
\end{equation}
with $v_{\mathrm{sw}}=\Vin$ while the high side is ON and $v_{\mathrm{sw}}=0$
while it is OFF. Define the ramp rates
\begin{equation}
\theta=-\frac{r}{L},\qquad
m_1=\frac{\Vin-\Vref}{L},\qquad
m_2=\frac{\Vref}{L},\qquad
m_\Sigma=m_1+m_2=\frac{\Vin}{L}.
\end{equation}

After a load step the controller works in the \emph{shifted state} measured at the
step (all in SI units), taken from the Pad\'e current/load estimator:
\begin{equation}
z_i=i_L-\hat I_{\mathrm{load,new}},\qquad
z_v=\Vout-\Vref,\qquad
\hat I=\hat I_{\mathrm{load,new}} .
\end{equation}
In the code (\texttt{buck\_optimal\_service}), $i_L=$ \texttt{cap\_iavg\_q20},
$\hat I=$ \texttt{cap\_iload\_q20} (both Q20 A), and
$\Vout=$ \texttt{cap\_adc}$\times$\texttt{UV\_PER\_COUNT}. The goal is to drive
$(z_i,z_v)\to(0,0)$ in \emph{one} ON--OFF action of durations $\Ton,\Toff$.

\section{Terminal equations}
Propagating the shifted state through the ON interval of length $x=\Ton$ then the
OFF interval of length $y=\Toff$ with the second--order bilinear (Pad\'e)
state--transition and requiring the final shifted state to vanish yields two
coupled equations (TPEL Sec.~IV, eqs.~35--36; verified against a direct Newton
solve to $\sim\!10^{-16}$ in \texttt{Tests/test\_optimal.c}):
\begin{align}
0&=z_i+p\,x+q\,y+g\,x\,y, \label{eq:e1}\\
0&=z_v+\tfrac{z_i}{2C}\,(x+y)+b\,x\,y, \label{eq:e2}
\end{align}
with coefficients
\begin{align}
p&=\tfrac{\theta}{2}z_i-\tfrac{z_v}{2L}+m_1-\theta\hat I, &
q&=\tfrac{\theta}{2}z_i-\tfrac{z_v}{2L}-m_2-\theta\hat I,\\
g&=\tfrac14\!\left(\theta^2-\tfrac{1}{LC}\right)z_i-\tfrac{\theta z_v}{4L}
   +\tfrac{\theta m_\Sigma}{2}, &
b&=\tfrac{\theta z_i}{4C}-\tfrac{z_v}{4LC}+\tfrac{m_\Sigma}{2C}.
\end{align}
A useful identity used below is
\begin{equation}
p-q=m_1+m_2=m_\Sigma=\frac{\Vin}{L}\ \Longrightarrow\ \frac{1}{p-q}=\frac{L}{\Vin}.
\label{eq:pmq}
\end{equation}

\section{Reduction to one quadratic (single square root)}
Let $s=x+y=\Ton+\Toff$ and $b_1=\dfrac{z_i}{2C}$.

\paragraph{Step 1 --- eliminate the bilinear term.}
From \eqref{eq:e2}, $b\,x y=-(z_v+b_1 s)$, i.e.
\begin{equation}
x\,y=-\frac{z_v+b_1 s}{b}. \label{eq:xy}
\end{equation}

\paragraph{Step 2 --- make $\Ton$ affine in $s$.}
Substitute $y=s-x$ and the product \eqref{eq:xy} into \eqref{eq:e1}. The bilinear
term $g\,x y$ becomes $-g\,(z_v+b_1 s)/b$, leaving a \emph{linear} equation in $x$:
\begin{equation}
z_i+q\,s+(p-q)\,x-\frac{g}{b}\,(z_v+b_1 s)=0 .
\end{equation}
Solving for $x$ and using \eqref{eq:pmq},
\begin{equation}
\boxed{\;x=\Ton=a_0+a_1 s,\qquad
a_0=-\frac{L}{\Vin}\!\left(z_i-\frac{g\,z_v}{b}\right),\quad
a_1=-\frac{L}{\Vin}\!\left(q-\frac{g\,b_1}{b}\right).}
\label{eq:affine}
\end{equation}

\paragraph{Step 3 --- one quadratic in $s$.}
With $x=a_0+a_1 s$ and $y=s-x=(1-a_1)s-a_0$, the product
$x y=a_1(1-a_1)s^2+a_0(1-2a_1)s-a_0^2$. Inserting into \eqref{eq:e2} gives
\begin{equation}
A\,s^2+B\,s+C_c=0,\qquad
\begin{aligned}
A&=b\,a_1(1-a_1),\\
B&=b\,a_0(1-2a_1)+b_1,\\
C_c&=z_v-b\,a_0^2,
\end{aligned}
\label{eq:quad}
\end{equation}
whose solution requires a \emph{single} square root:
\begin{equation}
s=\frac{-B\pm\sqrt{B^2-4A\,C_c}}{2A},\qquad
\Ton=a_0+a_1 s,\qquad \Toff=s-\Ton .
\label{eq:sol}
\end{equation}
The root ($+$ or $-$) is chosen as the physical one giving $\Ton>0$ and
$\Toff>0$. If the discriminant is negative, or $A=b=0$, no valid single--shot
solution exists and no pulse is issued.

\medskip
\noindent\textit{Note.} The paper's boxed closed form does not satisfy
\eqref{eq:e1}--\eqref{eq:e2}; \eqref{eq:affine}--\eqref{eq:sol} is the re--derived
form actually implemented, and it matches a Newton solve of
\eqref{eq:e1}--\eqref{eq:e2} to floating--point precision.

\section{As implemented: \texttt{buck\_optimal\_onoff()}}
\begin{lstlisting}[language=C]
q  = th*zi/2 - zv/(2L) - m2 - th*Ihat;
g  = 0.25*(th*th - 1/(L*C))*zi - th*zv/(4L) + th*(m1+m2)/2;
b  = th*zi/(4C) - zv/(4L*C) + (m1+m2)/(2C);      b1 = zi/(2C);
a0 = -(L/Vin)*(zi - g*zv/b);   a1 = -(L/Vin)*(q - g*b1/b);
A  = b*a1*(1-a1);   B = b*a0*(1-2*a1) + b1;   Cc = zv - b*a0*a0;
disc = B*B - 4*A*Cc;                 // < 0  -> no solution
s  = ((+/-)sqrt(disc) - B)/(2A);     // pick root with Ton>0, Toff>0
Ton = a0 + a1*s;   Toff = s - Ton;   // then clamp to [tmin, *_max]
\end{lstlisting}
This float form is now the \emph{host reference} (\texttt{Tests/test\_optimal.c}).
It ran in thread mode until September 2026, when the bench showed the pulse
arriving $\sim$100\,$\mu$s after the load edge: with the armed ISR at
$\sim$244 of 280 clocks the thread gets 15--35 clocks per period, and this
solve needs $\sim$50 periods of that. The firmware now solves in the ISR.

\section{ISR closed form (September 2026): two square roots, one cycle}\label{sec:isr}
For a load step--up write $a=i_0-I_{\rm eq}<0$ ($i_0$ the valley at the pulse
start, $I_{\rm eq}=\hat I+b$ the equilibrium current), $\Delta V=V_{\rm ref}-v_{\rm out}>0$,
$m_1=(V_{\rm in}-V_{\rm ref})/L$. The terminal state is the \emph{periodic steady
state} of the PWM converter at $I_{\rm eq}$: the valley $I_{\rm eq}-h$,
$h=\tfrac12(V_{\rm in}-V_{\rm ref})(V_{\rm ref}/V_{\rm in})T_{\rm sw}/L$ (half the
ripple, 2.45\,A here), so that the next period is an ordinary one and the first
sample after the pulse sits at the same point of the ripple as every steady--state
sample. The fall from $I_{\rm eq}+x$ down to $I_{\rm eq}-h$ sees the series drop,
\begin{equation}
m_2'=\frac{V_{\rm ref}+r\,i_2}{L},\qquad i_2=I_{\rm eq}+\frac{x-h}{2},
\end{equation}
($r$ = \texttt{g\_buck\_opt\_r\_off\_mohm}, a live tunable defaulting to the fitted
\texttt{BUCK\_R\_MOHM}). The pulse delivers, above $I_{\rm eq}$, the ON triangle
$(x^2-a^2)/(2m_1)$ plus the OFF area $(x^2-h^2)/(2m_2')$; equating the sum to the
deficit $C\,\Delta V$ gives, with $m_1+m_2'\simeq V_{\rm in}/L$,
\begin{equation}
x^2=\frac{V_{\rm ref}+r i_2}{V_{\rm in}}\Big[\frac{2C\,(V_{\rm in}-V_{\rm ref})}{L}\,\Delta V+a^2
+\frac{(V_{\rm in}-V_{\rm ref})\,h^2}{V_{\rm ref}+r i_2}\Big],
\qquad \Ton=\frac{x-a}{m_1},\qquad \Toff=\frac{(x+h)\,L}{V_{\rm ref}+r i_2}.
\end{equation}
The ideal--fall root $x_0$ ($r=0$, $h=0$) gives $i_2$, a second root gives $x$.
On the plant the model assumes (series drop on both ramps, trapezoidal average)
the pulse lands within 1.7\,mV and 2\,mA of $(V_{\rm ref},\,I_{\rm eq}-h)$
(\texttt{Tests/test\_solve.c}). The former ideal fall $\Toff=xL/V_{\rm ref}$ was up
to 2.4$\times$ too long at 10\,A with $r=114$\,m$\Omega$: on the bench every pulse
swung the current far below the load and the adaptation ran to its bounds. Applied as in Section~\ref{sec:applied}: the
excursion is capped so that $I_{\rm new}+x_{\rm app}\le I_{\rm pk}^{\max}$
(\texttt{sat}), the rise above $i_0$ is scaled by $\gamma_{\rm on}$ and $\Ton$ is
obtained from the resulting peak by the \emph{same} trapezoidal inversion the
current loop uses every cycle, $\Ton=L\,(I_{\rm pk}-i_L)/(V_L-r(i_L+I_{\rm pk})/2)$
(\texttt{buck\_cm\_ton\_raw}, one bounded UDIV), so the pulse peak lands where
the estimator predicts it, and $\Toff=\gamma_{\rm off}\,x_{\rm app}L/V_{\rm ref}$
(Section~\ref{sec:jac}); then the time clamps and the stretched--period floor of
(c)--(d). The terminal current reported to the reseed is the applied peak minus
what the final off--time removes, clamps and floor included.

\paragraph{Fixed point.} Currents in Q8\,A inside the solve ($a^2\le(32\,{\rm A})^2$
fits Q16\,A$^2$), $\Delta V$ clamped to 2047 counts so the first product stays below
$2^{31}$, the factors $V/V_{\rm in}$ applied as Q32 high--word multiplies
(\texttt{UMULL}), the $h^2$ term and $\Toff$ by bounded \texttt{UDIV}s. The roots are \texttt{buck\_isqrt32}: normalise $v\to v\,2^{2e}\in[2^{30},2^{32})$
with \texttt{CLZ}, read $\sqrt{i\,2^{24}}$ and its slope ($i=64..255$, one 32--bit word
each, 192 entries filled at start--up) and interpolate linearly, then shift back by
$e$: $\sim$18 clocks, within 1\,LSB of $\lfloor\sqrt v\rfloor$ over the 32--bit
range. The coefficients that depend on $V_{\rm ref}$ and $k$ are converted once by
the thread (\texttt{buck\_so\_coef\_from\_float}, struct \texttt{g\_buck.so}).

\paragraph{Timing (measured on the NUCLEO).} The ISR is entered 134--142
clocks into the switching period (ADC trigger at 600 ticks, trigger resync,
12.5-cycle conversion at 56\,MHz, EOC resync, exception entry), and the HRTIM
preload window closes at 264 clocks; the solve takes 172 clocks on silicon.
The pulse is therefore pipelined: trigger sample in period $n$ (end of the
hot path; it clears the hot path's \texttt{running} flag), solve in $n+1$ through
the existing cold entry (no HRTIM write, no deadline; the solve is done for
the state at the $n+3$ rollover: valley two periods ahead with the frozen
duty, droop plus $(I_{\rm new}-i_{\rm avg})\,2T/C$), \texttt{PER}/\texttt{CMP1}
written first thing in $n+2$ (measured 152--198 clocks into the period, 66--112
clocks before the deadline), stretched period $n+3$. Three switching periods
(5\,$\mu$s) from the trigger sample, against $\sim$53 with the thread solve.
The every--cycle path is unchanged (244 clocks, static count).

\section{Dual correction on the measured terminal state (manuscript Sec.~V)}\label{sec:jac}
The applied pulse is $\mathbf t_a=\boldsymbol\Gamma\,\mathbf t^\star$,
$\boldsymbol\Gamma=\mathrm{diag}(\gamma_{\rm on},\gamma_{\rm off})$ (eq.~40 of
the manuscript): $\gamma_{\rm on}$ scales the physical on--time, i.e.\ the rise
above the pulse--start current $i_0$, $\gamma_{\rm off}$ the off--time. The
closed form (Section~\ref{sec:isr}) aims at the periodic steady state
$(I_{\rm eq}-h,\,V_{\rm ref})$, $I_{\rm eq}=\hat I+b$ the load measured one period
after the edge plus a learned correction $b$.

Since the inductor current is sensed, the terminal state of the pulse is
\emph{observed}, at the first sample after the stretched period (600 ticks into
the first normal period, before the PI has changed anything):
\begin{align}
z_{v,f}&=v_s-V_{\rm ref}-R_{\rm ESR}\,(i_s-I_{\rm settled}),\\
z_{i,f}&=(i_s-m_1t_s)-(I_{\rm eq}-h),
\end{align}
where $v_s$, $i_s$ are the sampled voltage and current, $m_1t_s$ the rise from
the valley to the sample instant (0.96\,A at 61\,ns), and $I_{\rm eq}-h$ the target
valley. The equilibrium itself is checked when the linear--controller transition
ends (192 cycles, Section~\ref{sec:reseed}):
\begin{equation}
e_b=I_{\rm settled}-I_{\rm eq},
\end{equation}
$I_{\rm settled}$ the low--passed eq.~7 estimate at that moment ($\pm0.15$\,A,
against $\pm0.5$\,A for the former one--period window). Three residuals, three
parameters, one thread event per pulse. With the terminal map from the pulse
start,
\begin{align}
z_{i,f}&=z_{i0}+m_1\Ton-m_2\Toff,\\
z_{v,f}&=z_{v0}+\tfrac1C\big[z_{i0}(\Ton+\Toff)+\tfrac{m_1}{2}\Ton^2+m_1\Ton\Toff-\tfrac{m_2}{2}\Toff^2\big],
\end{align}
$m_2=(V_{\rm ref}+r\,i_2)/L$ the fall slope of this pulse ($i_2$ the mean current
of the fall), the Jacobian at the applied times is
\begin{equation}
\mathbf J=\begin{bmatrix} m_1 & -m_2\\[2pt]
\dfrac{z_{i0}+m_1(\Ton+\Toff)}{C} & \dfrac{z_{i0}+m_1\Ton-m_2\Toff}{C}\end{bmatrix},
\end{equation}
and the \emph{timing} correction is the manuscript's Newton step (eqs.~43--48)
on the full measured residual,
\begin{equation}
\boldsymbol\gamma^{+}=\Pi_\Gamma\!\left[\boldsymbol\gamma-\mu\,\mathrm{diag}(\mathbf t^\star)^{-1}\mathbf J^{-1}\begin{bmatrix}z_{i,f}\\ z_{v,f}\end{bmatrix}\right],
\qquad \mu=0.25,\quad \Gamma=[0.5,\,2.0],
\end{equation}
with at most $0.25$ of change per pulse and $\det\mathbf J$ floored at
$m_2\cdot0.5\,{\rm A}/C$. A clamped on--time (peak limit or $\Ton^{\max}$) cannot
move the delivered charge: $\Toff$ alone then takes $\Delta\Toff=\mu\,z_{i,f}/m_2$
(the voltage residual is unreachable). The \emph{equilibrium} correction is
\begin{equation}
b^{+}=b+\mu_b\,e_b,\qquad \mu_b=0.5,\quad |b|\le 3\,{\rm A}.
\end{equation}
The residuals are separable: a wrong equilibrium moves $e_b$ and not the
terminal residuals (the pulse still lands on its target), a wrong $L$, $C$ or $r$
moves the terminal residuals. \emph{Validity:} $z_{i,f}$, $z_{v,f}$, $e_b$ are
clipped to $\pm2$\,A / $\pm40$ counts before use, and a residual beyond $\pm6$\,A /
$\pm200$ counts, or a trigger on a droop beyond 200 counts, marks the record as
corrupted: that pulse teaches nothing. (On the bench, 25 September 2026, the
voltage sample read zero counts; the previous clip--only rule fired a pulse on
every unmasked cycle and ran both learners to their bounds.) Host test
(\texttt{Tests/test\_solve.c}, the learners fed with the plant's actual terminal
current and true load): a matched plant leaves $\gamma=1$, $b=0$; a 0.6\,A load
bias is found by $b$ with $\gamma$ untouched; $L$, $C$ and $r$ mismatches are
absorbed within 21 pulses (the contraction per pulse is $\mu$ times the
plant--to--model fall--slope ratio); a peak--clamped pulse has its current
residual nulled by $\Toff$ alone.

\paragraph{Post--pulse operating point (reseed).}\label{sec:reseed} With the
current measured the voltage loop's reference is real amperes: after the pulse it
is reseeded to $I_{\rm eq}$ itself and the current loop to the duty that holds it,
$d_{\rm trig}+rT_{\rm sw}(I_{\rm eq}-\bar i_{\rm trig})/V_{\rm in}$ (the fitted $r$ is
the DC duty--to--current relation), both precomputed in the write cycle. The
trigger is then masked for 192 cycles (320\,$\mu$s, the linear--controller
transition of manuscript Sec.~VI) while the PI owns the recovery; when the mask
ends the settled load closes the transient's record ($e_b$ above) and the thread
runs the dual correction. Without that mask a wrong reseed rings the loop, the
ring's undershoot re--triggers and the loop limit--cycles (simulation: 19 pulses
and a 1.2\,V overshoot on a 1$\to$6\,A step with a 15\,m$\Omega$ plant). Each
transient writes one of two alternating records (trigger sample, solve inputs
and outputs, reseed, terminal sample, settled load): no field the thread reads is
ever live.

\paragraph{Measured current (average current mode).} Since the fourth revision
the inductor current is sensed (1\,m$\Omega$ $\times$ 40 on A1, converted by ADC2 at the
same instant as the voltage): the inner loop is an average--current PI, the load
estimate is eq.~7 of the manuscript on the measured period average
($\bar i_L=i_{\rm sample}+(V_{\rm in}-V_{\rm ref})/L\,(\Ton/2-t_s)$, $t_s$ the sample's
position on the on--ramp), $i_0$ is the measured valley extrapolated two periods
with the measured drift and $\hat I$ the load measured over the period after the
edge. The Pad\'e recursion no longer runs in the ISR.

\paragraph{Capacitor ESR.} The ADC sample is $v_C+R_{\rm ESR}\,(i_L-I_{\rm load})$ at
the sample instant (on the bench: a $+80$--$100$\,mV spike during the pulse at
$\sim$7.6\,A of capacitor current, i.e.\ 10--15\,m$\Omega$). Two corrections use the
\emph{measured} sample current: the solve refills
$\Delta V=V_{\rm ref}-v_s+R_{\rm ESR}\,(i_s-\hat I)$ rather than the raw error
(during a droop the inductor lags the load, so the sample sits below the
capacitor by that much: 30 counts for a 2\,A lag at 12\,m$\Omega$, more than half
the droop of a 3\,A step), and the terminal--voltage residual is corrected against
the settled load as above. $R_{\rm ESR}$ is a live tunable
(\texttt{g\_buck\_opt\_esr\_mohm}) measured on the scope as the spike during the
pulse divided by $I_{\rm pk}-I_{\rm load}$; the former in--loop identification read
60--75\,\% of the truth and is gone.

\section{Applied $\Ton/\Toff$: gain, peak--current limit, clamps}\label{sec:applied}
The solver output $(\Ton^\star,\Toff^\star)$ is post--processed in
\texttt{buck\_optimal\_service()} before being written to the HRTIM:

\subsection*{(a) Timing gains}
The excursion $x$ of the closed form is capped at the peak limit
($x_{\rm app}\le I_{\rm pk}^{\max}-I_{\rm new}$, flag \texttt{sat}); the applied rise
above $i_0$ is $\gamma_{\rm on}(I_{\rm new}+x_{\rm app}-i_0)$, re--capped at the
peak limit, and the on--time is its trapezoidal inversion,
\begin{equation}
\Ton'=\frac{L\,(I_{\rm pk,a}-i_0)}{V_L-r\,(i_0+I_{\rm pk,a})/2},\qquad
\Toff'=\gamma_{\rm off}\,\frac{x_{\rm app}}{m_2},
\end{equation}
with $\gamma_{\rm on},\gamma_{\rm off}$ from Section~\ref{sec:jac} (1 at boot).
The former learned gain $k$ (excursion scaling with an empirical secant on the
first post--pulse sample) is gone: its residual depended on the trigger timing,
not on $k$.

\subsection*{(b) Peak--current--mode limit (control, not protection)}
The on--time is capped so the predicted peak inductor current stays at a
configured limit $\Ipk$ (\texttt{BUCK\_OPT\_IPK\_MAX\_A}). With the estimator's
inductor current $i_L$ at the step as the ramp start,
\begin{equation}
I_{\mathrm{peak}}=i_L+\frac{\Vin-\Vout}{L}\,\Ton\le \Ipk
\;\Longrightarrow\;
\Ton\le \Ton^{\Ipk}=\frac{(\Ipk-i_L)\,L}{\Vin-\Vout}.
\end{equation}
This is the normally binding limit; it only shortens $\Ton$ and never trips a
fault (over--current \emph{protection} is external to the loop).

\subsection*{(c) Time clamp and quantization}
Finally
\begin{equation}
\Ton^{\mathrm{app}}=\min\!\bigl(\Ton',\,\Ton^{\Ipk},\,\Ton^{\max}\bigr),\ \ge t_{\min},
\qquad
\Toff^{\mathrm{app}}=\max\!\bigl(\Toff',\,t_{\min}\bigr),
\end{equation}
where $\Ton^{\max}<\Tsw$ is required so that the compare write $\texttt{CMP1}=\Ton$
is a valid compare against the (still active) one--period register. If either the
peak or time clamp binds, a saturation flag excludes $\gamma_{\rm on}$ from the
adaptation of that pulse. The
seconds values are quantized to HRTIM ticks,
\begin{equation}
n_{\mathrm{on}}=\Bigl\lfloor \Ton^{\mathrm{app}}/T_{\mathrm{tick}}+\tfrac12\Bigr\rfloor,\qquad
n_{\mathrm{off}}=\Bigl\lfloor \Toff^{\mathrm{app}}/T_{\mathrm{tick}}+\tfrac12\Bigr\rfloor,
\end{equation}
and applied for one stretched period as $\texttt{PER}=n_{\mathrm{on}}+n_{\mathrm{off}}$,
$\texttt{CMP1}=n_{\mathrm{on}}$ (high side ON for $n_{\mathrm{on}}$, OFF for the
remainder).

\subsection*{(d) Stretched--period floor}
The ISR that restores the normal PWM runs at a \emph{fixed} point of the
stretched period (ADC trigger at \texttt{CMP3} plus conversion and interrupt
entry, roughly $4.5$--$5.3\times10^3$ ticks in). A stretched period shorter than
that would end before the restore is written, the short pulse would repeat and
the deadline check would trip. The firmware therefore enforces
\begin{equation}
n_{\mathrm{on}}+n_{\mathrm{off}}\;\ge\;N_{\min}
=\texttt{BUCK\_OPT\_PER\_MIN\_TICKS}=\texttt{BUCK\_HRTIM\_PERIOD},
\qquad n_{\mathrm{off}}\leftarrow N_{\min}-n_{\mathrm{on}}\ \text{if needed},
\end{equation}
i.e.\ the large--signal period is never shorter than a normal one, so every ISR
timing assumption of a normal cycle holds inside it. With $k=1$ the solver's
$\Ton^\star+\Toff^\star$ is $2$--$4\,\mu$s for realistic captured states and the
floor does not bind; it binds only when $\gamma_{\rm off}$ has shrunk $\Toff$.

\subsection*{(e) Trigger}
The state--anchored trigger that starts a transient acts on the output--capacitor
state, both inputs in ADC counts: $z_v=v_s-V_{\rm ref}$ and $d_v=v_s(n)-v_s(n-1)$,
the one--period change ($C\,dv/dt$ is the capacitor current $i_L-I_{\rm load}$: the
load step itself, before the current loop has moved the inductor; 1\,A is 8 counts
per period). With the inductor current regulated, the former current--displacement
axis only reported the current loop. Fire when
$S=(z_v/V_a)^2+(d_v/dV_a)^2+2\rho\,(z_v/V_a)(d_v/dV_a)>1$ with $z_v<0$ and $d_v\le0$
(below the reference and not recovering: the only signature the step--up solver
can act on). $S$ is evaluated inside the ISR in integer arithmetic by completing
the square, $S=u^2+w^2$ with $u=n_v+\rho n_d$, $w=\sqrt{1-\rho^2}\,n_d$ ($n_v=z_v/V_a$,
$n_d=d_v/dV_a$ in Q14, compared against $2^{28}$), only for samples that pass the
two sign tests. $V_a$ (25 counts), $dV_a$ (16 counts per period), $\rho$ remain
run--time float tunables; the thread converts them into the ISR's integer
coefficients whenever they change. A fast 3\,A step scores 2.25 in its first period;
a 15--40\,$\mu$s slew fires on the level, two thirds of the way down the dip the
current--mode PI leaves on its own; $\pm6$ counts of noise on both samples score
0.6 (simulation: $\pm8$ counts false--trigger on the slope axis, so $dV_a$ is set
from the steady--state score of the candidate samples, which the panel shows). A feedback sample that jumps more than 256 counts in one period is
held off for that cycle (two consecutive such samples are accepted), and a sample
below a quarter of the reference for 500\,$\mu$s trips the feedback--lost fault.

\paragraph{Model scale and dead time.} $V_{\rm ref}$ and every voltage derived
from the ADC use the \emph{measured} scale, 793\,$\mu$V per count (scope mean
792.6\,mV at the 1000--count setpoint; an assumed 1\,mV per count made every
model voltage 26\,\% too large, so the pulse refilled a 26\,\% larger deficit
with a 26\,\% too steep off--ramp and ended $\sim$1.3\,A above the load). The
inductor sees the CMP1 on--time minus the 50\,ns rising dead time (272 ticks) in
the estimator table, the current--loop inversion and the pulse solve.

\section{Parameters}
\begin{center}
\begin{tabular}{@{}lll@{}}
\toprule
Symbol & Macro & Value \\
\midrule
$\Vin$ & \texttt{BUCK\_VIN\_V} & $6.0$ V\\
$\Vref$ & \texttt{BUCK\_VOUT\_MV} & $0.793$ V (measured)\\
$L$ & \texttt{BUCK\_L\_H} & $330$ nH\\
$C$ & \texttt{BUCK\_C\_F} & $260\,\mu$F\\
$r$ & \texttt{BUCK\_R\_MOHM} & $114$ m$\Omega$\\
$t_{\min}$ & \texttt{BUCK\_OPT\_TMIN\_S} & $100$ ns\\
$\Ton^{\max}$ & \texttt{BUCK\_OPT\_TON\_MAX\_S} & $1.5\,\mu$s ($<\Tsw=1.667\,\mu$s)\\
$\Toff^{\max}$ & \texttt{BUCK\_OPT\_TOFF\_MAX\_S} & $6\,\mu$s\\
$\Ipk$ & \texttt{BUCK\_OPT\_IPK\_MAX\_A} & $10$ A\\
$\mu$, $\Gamma$ & \texttt{BUCK\_OPT\_GAMMA\_*} & $0.25$, $[0.5,2.0]$\\
$\mu_b$, $|b|_{\max}$ & \texttt{BUCK\_OPT\_IEQ\_*} & $0.5$, $3$ A\\
$T_{\mathrm{tick}}$ & $\Tsw/\text{PERIOD}$ & $1.667\,\mu\text{s}/8960$\\
\bottomrule
\end{tabular}
\end{center}

\end{document}
